% Options for packages loaded elsewhere
\PassOptionsToPackage{unicode}{hyperref}
\PassOptionsToPackage{hyphens}{url}
\PassOptionsToPackage{dvipsnames,svgnames,x11names}{xcolor}
\documentclass[
  spanish,
  11pt,
]{article}
\usepackage{xcolor}
\usepackage[margin=2.5cm]{geometry}
\usepackage{amsmath,amssymb}
\setcounter{secnumdepth}{-\maxdimen} % remove section numbering
\usepackage{iftex}
\ifPDFTeX
  \usepackage[T1]{fontenc}
  \usepackage[utf8]{inputenc}
  \usepackage{textcomp} % provide euro and other symbols
\else % if luatex or xetex
  \usepackage{unicode-math} % this also loads fontspec
  \defaultfontfeatures{Scale=MatchLowercase}
  \defaultfontfeatures[\rmfamily]{Ligatures=TeX,Scale=1}
\fi
\usepackage{lmodern}
\ifPDFTeX\else
  % xetex/luatex font selection
\fi
% Use upquote if available, for straight quotes in verbatim environments
\IfFileExists{upquote.sty}{\usepackage{upquote}}{}
\IfFileExists{microtype.sty}{% use microtype if available
  \usepackage[]{microtype}
  \UseMicrotypeSet[protrusion]{basicmath} % disable protrusion for tt fonts
}{}
\makeatletter
\@ifundefined{KOMAClassName}{% if non-KOMA class
  \IfFileExists{parskip.sty}{%
    \usepackage{parskip}
  }{% else
    \setlength{\parindent}{0pt}
    \setlength{\parskip}{6pt plus 2pt minus 1pt}}
}{% if KOMA class
  \KOMAoptions{parskip=half}}
\makeatother
\usepackage{longtable,booktabs,array}
\newcounter{none} % for unnumbered tables
\usepackage{calc} % for calculating minipage widths
% Correct order of tables after \paragraph or \subparagraph
\usepackage{etoolbox}
\makeatletter
\patchcmd\longtable{\par}{\if@noskipsec\mbox{}\fi\par}{}{}
\makeatother
% Allow footnotes in longtable head/foot
\IfFileExists{footnotehyper.sty}{\usepackage{footnotehyper}}{\usepackage{footnote}}
\makesavenoteenv{longtable}
\ifLuaTeX
\usepackage[bidi=basic,shorthands=off]{babel}
\else
\usepackage[bidi=default,shorthands=off]{babel}
\fi
\ifLuaTeX
  \usepackage{selnolig} % disable illegal ligatures
\fi
\setlength{\emergencystretch}{3em} % prevent overfull lines
\providecommand{\tightlist}{%
  \setlength{\itemsep}{0pt}\setlength{\parskip}{0pt}}
% Símbolos Unicode (griegos, sub/superíndices, operadores) que Latin Modern no trae
\usepackage{newunicodechar}
\newfontfamily\simbolos{DejaVu Serif}[Scale=MatchLowercase]
\newunicodechar{α}{{\simbolos α}}
\newunicodechar{β}{{\simbolos β}}
\newunicodechar{δ}{{\simbolos δ}}
\newunicodechar{η}{{\simbolos η}}
\newunicodechar{π}{{\simbolos π}}
\newunicodechar{ρ}{{\simbolos ρ}}
\newunicodechar{τ}{{\simbolos τ}}
\newunicodechar{χ}{{\simbolos χ}}
\newunicodechar{ᵢ}{{\simbolos ᵢ}}
\newunicodechar{⁰}{{\simbolos ⁰}}
\newunicodechar{⁶}{{\simbolos ⁶}}
\newunicodechar{⁷}{{\simbolos ⁷}}
\newunicodechar{⁹}{{\simbolos ⁹}}
\newunicodechar{⁻}{{\simbolos ⁻}}
\newunicodechar{ⁿ}{{\simbolos ⁿ}}
\newunicodechar{₀}{{\simbolos ₀}}
\newunicodechar{₁}{{\simbolos ₁}}
\newunicodechar{₂}{{\simbolos ₂}}
\newunicodechar{₊}{{\simbolos ₊}}
\newunicodechar{ₖ}{{\simbolos ₖ}}
\newunicodechar{ₙ}{{\simbolos ₙ}}
\newunicodechar{∈}{{\simbolos ∈}}
\newunicodechar{∉}{{\simbolos ∉}}
\newunicodechar{∝}{{\simbolos ∝}}
\newunicodechar{≈}{{\simbolos ≈}}
\newunicodechar{≠}{{\simbolos ≠}}
\newunicodechar{≤}{{\simbolos ≤}}
\newunicodechar{≥}{{\simbolos ≥}}
\newunicodechar{⊆}{{\simbolos ⊆}}
\usepackage{bookmark}
\IfFileExists{xurl.sty}{\usepackage{xurl}}{} % add URL line breaks if available
\urlstyle{same}
\hypersetup{
  pdftitle={Algoritmos Heurísticos: Ciclo Hamiltoniano sobre el FHCP Challenge Set},
  pdflang={es},
  colorlinks=true,
  linkcolor={Maroon},
  filecolor={Maroon},
  citecolor={Blue},
  urlcolor={Blue},
  pdfcreator={LaTeX via pandoc}}

\title{Algoritmos Heurísticos: Ciclo Hamiltoniano sobre el FHCP
Challenge Set}
\author{}
\date{}

\begin{document}
\maketitle

Código de la práctica ``Aula invertida -- Algoritmos heurísticos'':
carga grafos del FHCP, ejecuta metaheurísticas (GA, SA, ACO, Bee Colony,
DE, PSO) y devuelve un arreglo de vértices (el ciclo propuesto).

\section{Estructura}\label{estructura}

\begin{itemize}
\tightlist
\item
  \texttt{src/core/} loader HCP, lector/escritor TOUR, fitness,
  validador, representación, operadores, interfaz común
\item
  \texttt{src/algorithms/} \texttt{ga.py}, \texttt{sa.py} (swap y
  2-opt), \texttt{aco.py}, \texttt{bee.py}, \texttt{de.py},
  \texttt{pso.py}, \texttt{random\_search.py} + registro
\item
  \texttt{src/core/preprocess.py} aristas forzadas de vértices de grado
  2; \texttt{operators.py} también trae ER (recombinación de aristas)
\item
  \texttt{experiments/} \texttt{runner.py} (una corrida),
  \texttt{benchmark.py} (todo el benchmark), \texttt{tuning.py} (ajuste
  de parámetros), \texttt{analisis.py} (tablas, pruebas estadísticas y
  figuras), \texttt{catalogo.py} (catálogo de instancias),
  \texttt{interfaz.py} (interfaz gráfica con tkinter)
\item
  \texttt{data/instances/} instancias (.hcp); \texttt{data/README.md}
  explica su origen; \texttt{data/catalogo\_fhcp.csv} cataloga las 1001
\item
  \texttt{INFORME.md} informe técnico completo (teoría, comparativa,
  benchmarking contra el FHCP y conclusiones)
\item
  \texttt{results/} resultados del benchmark, del ajuste, tablas y
  figuras; \texttt{tests/} pruebas (\texttt{pytest})
\end{itemize}

\section{Convenciones}\label{convenciones}

\begin{itemize}
\tightlist
\item
  Vértices internos 0..n-1; los archivos usan 1..n (la conversión está
  en \texttt{hcp\_io.py}).
\item
  Fitness = aristas faltantes del ciclo cerrado; éxito si fitness == 0.
  El validador es independiente de los algoritmos.
\item
  Todo algoritmo expone
  \texttt{run(instance,\ config,\ seed)\ -\textgreater{}\ Result(best\_route,\ best\_fitness,\ history,\ time,\ iterations)}.
\item
  1 iteración = 1 evaluación de fitness (completa en GA/ACO/DE/PSO; por
  delta O(1) en SA/Bee). Presupuesto = evaluaciones o tiempo.
\end{itemize}

\section{Uso (desde la carpeta raíz del
proyecto)}\label{uso-desde-la-carpeta-rauxedz-del-proyecto}

\begin{verbatim}
pip install -r requirements.txt
python -m pytest

# Interfaz gráfica (tkinter): elegir grafo y algoritmo, ver el ciclo y la convergencia
python -m experiments.interfaz
# En Linux, si falta tkinter: sudo apt install python3-tk

# Una corrida con varias semillas -> CSV
python -m experiments.runner --instance data/instances/graph1.hcp --algorithm sa_2opt --seeds 10 --max-iter 100000 --out results/sa_graph1.csv

# Benchmark completo (reanudable) + análisis
python -m experiments.benchmark --instances data/instances/graph1.hcp data/instances/graph24.hcp \
    --algorithms ga sa sa_2opt aco bee de pso random_search --seeds 10 --max-evals 100000 --time-limit 15 \
    --out results/bench --tours results/tours
python -m experiments.analisis --prefix results/bench --out results/analisis
\end{verbatim}

Algoritmos disponibles: \texttt{ga}, \texttt{sa} (swap),
\texttt{sa\_2opt}, \texttt{aco}, \texttt{bee}, \texttt{de},
\texttt{pso}, \texttt{random\_search} y las mejoras de la segunda ronda:
\texttt{aco\_ls} (ACO + 2-opt, el único que resolvió una instancia),
\texttt{aco\_ls\_pre} (+ preprocesamiento) y \texttt{ga\_er} (GA con
recombinación de aristas). Los tours válidos encontrados se guardan en
\texttt{results/tours/} (formato TOUR del FHCP). Los parámetros por
algoritmo están en el docstring de cada módulo y se pasan en
\texttt{Config.params}.

\section{Autores}\label{autores}

Universidad Politécnica de Victoria · Docente: Dr.~Said Polanco Martagón

{\def\LTcaptype{none} % do not increment counter
\begin{longtable}[]{@{}ll@{}}
\toprule\noalign{}
Nombre & Matrícula \\
\midrule\noalign{}
\endhead
\bottomrule\noalign{}
\endlastfoot
Jared de Jesús Olazarán López & 2430022 \\
Georgina Reta Limas & 2430043 \\
Dulce Yaretzi Martínez Pérez & 2430038 \\
Melissa Jazmín Torres Martínez & 2430017 \\
Hugo Guillermo De León Ruiz & 2330419 \\
\end{longtable}
}

\end{document}
